\documentclass[10pt,a4paper]{amsart}
\usepackage[margin=19mm]{geometry}
\usepackage{amsmath,amssymb,mathtools}
\usepackage[scr=rsfs,cal=euler]{mathalfa}
\usepackage{xfrac}
\usepackage[unicode,hidelinks,bookmarks=false]{hyperref}
\newcommand{\ie}{i.e.\ }
\newcommand{\cf}{cf.\ }
\newcommand{\resp}{respectively\ }
\newcommand{\Subsection}{\subsection}
\providecommand{\nobreakdash}{\nobreak\hbox{-}}
\setlength{\emergencystretch}{2em}
\multlinegap0pt
\usepackage{bm}
\usepackage{bbm}
\usepackage{dsfont}
\usepackage{xspace}
\usepackage{cancel}
\usepackage{mathtools}
\usepackage{amsthm}
\makeatletter
\@ifundefined{Thm}{\newtheorem{Thm}{Theorem}}{}
\makeatother
\newcommand{\bbN}{\mathbb{N}}      % natural numbers %                     
\newcommand{\cF}{\mathcal{F}}
\newcommand{\lef}{L}
\newcommand{\rig}{R}
\title[Elastic rigid rod in an expanding universe]{Elastic rigid rod in an expanding universe}
\author{Sim\~ao Correia, Jos\'e Nat\'ario, Jorge Drumond Silva}
\date{Source: arXiv:2507.21227v1 (CC BY 4.0)}
\begin{document}
\maketitle

\noindent\emph{Source and licence.} Elastic rigid rod in an expanding universe. Version arXiv:2507.21227v1 (CC BY 4.0), distributed under the Creative Commons Attribution 4.0 International licence, as recorded in the arXiv licence metadata for this exact version and shown on the abstract page https://arxiv.org/abs/2507.21227v1. Full rights notice: licenses/arxiv-2507.21227-CC-BY-4.0.txt.\par\smallskip

\noindent\textit{Editorial scope.} Counted unit: the Thm whose source label is \texttt{GWPBA} in the article (source0.tex lines 856--858, source label \texttt{GWPBA}), delivered together with the article's own complete original proof of it (source0.tex lines 860--1045, one \texttt{proof} environment of 186 lines). No mathematics is written, repaired or re-derived: statement and proof are the article's own text line for line. Required context retained uncounted: IBVP (lines 220--226); L+- (lines 319--321); IBVP2 (lines 664--671). Every definition, display and label of those lines is retained unchanged. Omitted from this package and not redistributed: 1--219, 227--318, 322--663, 672--855, 859--859, 1046--1127. Nothing inside a retained excerpt is omitted or altered: every retained body is delivered exactly as the article's lines are written, so no omission receipt of any kind accompanies this package. Citations: the retained excerpts carry no citation command at all, so no bibliography entry of the article is transcribed and no reference is redistributed. Reference adaptation: none was needed and none was made; every reference inside the delivered unit resolves inside it, so no unresolved reference or citation is produced. Printed slips kept literal: no printed word, symbol, hypothesis or number of the counted statement or of its proof was altered. Layout and class: the article's own class is \texttt{amsart} with packages including amssymb, enumerate, epsf, psfrag, graphics, graphicx, xcolor, hyperref, mathrsfs, wrapfig, float, mathtools, caption, subfig; the wrapper is the lane's pinned \texttt{amsart} editorial wrapper at 10pt on A4, which restates the theorem-like environments the retained text needs and adds the article's own macro definitions that the retained text needs (\texttt{\textbackslash bbN}, \texttt{\textbackslash cF}, \texttt{\textbackslash lef}, \texttt{\textbackslash rig}), with extra wrapper packages (bm, bbm, dsfont, xspace, cancel, mathtools). The delivered numbering is the unit's own. Assets: no retained span contains a figure environment, a graphics-inclusion command, a PDF-page inclusion, a table or a TikZ picture; no image file is copied into this package, the package declares no asset dependency and no image file is redistributed with it. Licence: version arXiv:2507.21227v1 of this article is distributed under the Creative Commons Attribution 4.0 International licence, as recorded in the arXiv licence metadata for this exact version and shown on the abstract page https://arxiv.org/abs/2507.21227v1. A full rights notice is delivered at licenses/arxiv-2507.21227-CC-BY-4.0.txt. Characters: every retained excerpt is pure ASCII, so the delivered engine, XeLaTeX with its default Latin Modern text font, reports no missing character.
\begin{equation} \label{IBVP}
\begin{cases}
- \ddot{u} + u'' = K e^{- 2u} \\
u(0,\lambda)=\dot{u}(0,\lambda)=0, \quad \lambda \in [0,1] \\
u(\tau,0) = u(\tau,1) = 0, \quad \tau\geq 0
\end{cases} \, ,
\end{equation}

\begin{equation} \label{L+-}
L^- = \frac{\partial}{\partial \tau} - \frac{\partial}{\partial \lambda} \qquad \text{ and } \qquad L^+ = \frac{\partial}{\partial \tau} + \frac{\partial}{\partial \lambda} \, ,
\end{equation}

\begin{equation} \label{IBVP2}
\begin{cases}
- \ddot{u} + u'' = K e^{- 2u} \\
u(0,\lambda)=\dot{u}(0,\lambda)=0, \quad \quad \lambda \in [0,1] \\
u'(\tau,0) = \sigma \sinh u(\tau,0), \quad \,\,\,\, \tau\geq 0 \\
u'(\tau,1) = - \sigma \sinh u(\tau,1), \quad \tau\geq 0
\end{cases} \, ,
\end{equation}

\begin{Thm} \label{GWPBA}
The IBVPs~\eqref{IBVP} and \eqref{IBVP2} admit unique $C^2$ solutions on maximal time intervals $[0,T)$, where either $T = + \infty$ or $\limsup_{t {\scriptscriptstyle \nearrow} T} \| u(\tau, \cdot)\|_{L^\infty} = + \infty$.
\end{Thm}

\begin{proof}
We present only the proof for the more complicated case of the IBVP~\eqref{IBVP2}; the proof for the simpler IBVP~\eqref{IBVP} is a trivial adaptation.

Let us consider first the IBVP
\begin{equation} \label{IBVP3}
\begin{cases}
- \ddot{u} + u'' = K e^{- 2u} \\
u(\tau_0,\lambda)=f(\lambda), \quad \quad \lambda \in [0,1] \\
\dot{u}(\tau_0,\lambda)=g(\lambda), \quad \quad \lambda \in [0,1] \\
u'(\tau,0) = \sigma \sinh u(\tau,0), \quad \,\,\,\, \tau\in [\tau_0, \tau_1] \\
u'(\tau,1) = - \sigma \sinh u(\tau,1), \quad \tau\in [\tau_0, \tau_1]
\end{cases} \, ,
\end{equation}
where we assume that $\Delta \tau = \tau_1 - \tau_0 < \frac12$ and that $f \in C^2([0,1])$ and $g \in C^1([0,1])$ satisfy the compatibility conditions
\begin{equation} \label{compatibility1}
\begin{cases}
f'(0)=\sigma \sinh f(0) \\
f'(1)=-\sigma \sinh f(1)
\end{cases}
\end{equation}
and
\begin{equation} \label{compatibility2}
\begin{cases}
g'(0)=\sigma g(0) \cosh f(0) \\
g'(1)=-\sigma g(1) \cosh f(1)
\end{cases} \, .
\end{equation}
To solve this problem, we introduce the new variables
\begin{equation}
v = L^+ u \, , \qquad w = L^- u \, ,
\end{equation}
where $L^+$ and $L^-$ are the vector fields defined in \eqref{L+-}, and study the first order IBVP
\begin{equation} \label{IBVPfirstorder}
\begin{cases}
\dot{u} = \frac12 (v + w) \\
L^- v = - K e^{- 2u} \\
L^+ w = - K e^{- 2u} \\
u(\tau_0,\lambda)=f(\lambda), \quad \quad \lambda \in [0,1] \\
v(\tau_0,\lambda)=g(\lambda) + f'(\lambda), \quad \quad \lambda \in [0,1] \\
w(\tau_0,\lambda)=g(\lambda) - f'(\lambda), \quad \quad \lambda \in [0,1] \\
v(\tau,0) - w(\tau,0) = 2\sigma \sinh u(\tau,0), \quad \,\,\,\, \tau\in [\tau_0, \tau_1] \\
w(\tau,1) - v(\tau,1) = 2\sigma \sinh u(\tau,1), \quad \tau\in [\tau_0, \tau_1]
\end{cases} \, ,
\end{equation}which we write as the following system of integral equations:
\begin{equation} \label{IBVP4}
\begin{cases}
u(\tau,\lambda) = f(\lambda) + \frac12 \int_{\tau_0}^\tau \left[ v(s,\lambda) + w(s,\lambda) \right] ds \\
v(\tau,\lambda)=g(\lambda + \tau - \tau_0) + f'(\lambda + \tau - \tau_0) - \int_{\tau_0}^\tau K(s,\lambda+\tau-s) e^{- 2u(s,\lambda+\tau-s)} ds \, , \quad \tau + \lambda \leq 1 + \tau_0 \\
v(\tau,\lambda)= w_\rig(\tau, \lambda) - 2 \sigma \sinh u_\rig(\tau,\lambda) - \int_{\tau + \lambda - 1}^\tau K(s,\lambda+\tau-s) e^{- 2u(s,\lambda+\tau-s)} ds \, , \quad \tau + \lambda\geq 1 + \tau_0 \\
w(\tau,\lambda)=g(\lambda - \tau + \tau_0) - f'(\lambda - \tau + \tau_0) - \int_{\tau_0}^\tau K(s,\lambda-\tau+s) e^{- 2u(s,\lambda-\tau+s)} ds \, , \quad \tau - \lambda \leq \tau_0 \\
w(\tau,\lambda)=v_\lef(\tau,\lambda) - 2 \sigma \sinh u_\lef(\tau,\lambda) - \int_{\tau - \lambda}^\tau K(s,\lambda-\tau+s) e^{- 2u(s,\lambda-\tau+s)} ds \, , \quad \tau - \lambda \geq \tau_0 \\
u_\lef(\tau,\lambda) = f(0) + \frac12 \int_{\tau_0}^{\tau-\lambda} \left[ v(s,0) + w(s,0) \right] ds \\
u_\rig(\tau,\lambda) = f(1) + \frac12 \int_{\tau_0}^{\tau+\lambda-1} \left[ v(s,1) + w(s,1) \right] ds \\
v_\lef(\tau,\lambda)=g(\tau - \lambda - \tau_0) + f'(\tau - \lambda - \tau_0) - \int_{\tau_0}^{\tau-\lambda} K(s,\tau-\lambda-s) e^{- 2u(s,\tau-\lambda-s)} ds \\
w_\rig(\tau,\lambda)=g(2 - \tau - \lambda + \tau_0) - f'(2 - \tau - \lambda + \tau_0) - \int_{\tau_0}^{\tau+\lambda-1} K(s,2 - \tau - \lambda+s) e^{- 2u(s,2 - \tau - \lambda+s)} ds
\end{cases} \, .
\end{equation}

Here the first equation states that $u$ is the time integral of its time derivative, and the next four equations give $u$ and $v$ by integrating the equation along the characteristics, starting either at the initial time or at the boundaries; the functions $u_\lef$, $u_\rig$, $v_\lef$ and $w_\rig$ are the boundary values needed to start the integration at the boundaries, and are themselves obtained by integrations that always start at the initial time (since $\Delta \tau < \frac12$).

Consider the Banach space $X= \left(C^0\left([\tau_0, \tau_1] \times [0,1]\right)\right)^3$ with the norm
\begin{equation}
\|(u,v,w)\| = \| u \|_{L^\infty} + \| v \|_{L^\infty} + \| w \|_{L^\infty} \, ,
\end{equation}
and let $\cF:X \to X$ be the obvious operator implicit in~\eqref{IBVP4}, whose fixed points are the solutions of \eqref{IBVPfirstorder} (note that the compatibility conditions~\eqref{compatibility1} ensure that $\cF$ indeed maps continuous functions to continuous functions). If we write
\begin{equation}
\cF(u,v,w) = (\bar{u},\bar{v},\bar{w}) \, ,
\end{equation}
and assume that
\begin{equation}
\|(u,v,w)\| \leq R \, ,
\end{equation}
for a suitable constant $R>0$, then we have
\begin{equation}
\| \bar{u} \|_{L^\infty} \leq \| f \|_{L^\infty} + R \Delta \tau \, .
\end{equation}
Similarly, we have, for $\tau + \lambda \leq 1 + \tau_0$,
\begin{equation}
| \bar{v} | \leq \| g \|_{L^\infty} + \| f' \|_{L^\infty} + \| K \|_{L^\infty} e^{2R} \Delta \tau \, ,
\end{equation}
and, for $ \tau + \lambda\geq 1 + \tau_0$,
\begin{equation}
| \bar{v} | \leq \| w_\rig \|_{L^\infty} + 2 \sigma \| \sinh u_\rig \|_{L^\infty} + \| K \|_{L^\infty} e^{2R} \Delta \tau \, .
\end{equation}
Since
\begin{equation}
\| w_\rig \|_{L^\infty} \leq \| g \|_{L^\infty} + \| f' \|_{L^\infty} + \| K \|_{L^\infty} e^{2R} \Delta \tau \, ,
\end{equation}
and
\begin{equation}
\| u_\rig \|_{L^\infty} \leq \| f \|_{L^\infty} + R \Delta \tau \, ,
\end{equation}
we conclude that
\begin{equation}
\| \bar{v} \|_{L^\infty} \leq \| g \|_{L^\infty} + \| f' \|_{L^\infty} + 2 \| K \|_{L^\infty} e^{2R} \Delta \tau + 2 \sigma \sinh \left( \| f \|_{L^\infty} + R \Delta \tau \right)  \, ,
\end{equation}
and similarly
\begin{equation}
\| \bar{w} \|_{L^\infty} \leq \| g \|_{L^\infty} + \| f' \|_{L^\infty} + 2 \| K \|_{L^\infty} e^{2R} \Delta \tau + 2 \sigma \sinh \left( \| f \|_{L^\infty} + R \Delta \tau \right)  \, .
\end{equation}
We then have
\begin{equation}
\| (\bar{u},\bar{v},\bar{w}) \| \leq  \| f \|_{L^\infty} + 2\| g \|_{L^\infty} + 2\| f' \|_{L^\infty} + R \Delta \tau  + 4 \| K \|_{L^\infty} e^{2R} \Delta \tau + 4 \sigma \sinh \left( \| f \|_{L^\infty} + R \Delta \tau \right)  \, ,
\end{equation}
and so if we choose
\begin{equation} \label{sizeofR}
R =  \| f \|_{L^\infty} + 2\| g \|_{L^\infty} + 2\| f' \|_{L^\infty} + 4 \sigma \sinh \left( \| f \|_{L^\infty} \right) + 1 \, ,
\end{equation}
and then $\Delta \tau$ sufficiently small so that
\begin{equation} \label{sizeofdeltatau}
R \Delta \tau  + 4 \| K \|_{L^\infty} e^{2R} \Delta \tau + 4 \sigma \sinh \left( \| f \|_{L^\infty} + R \Delta \tau \right) - 4 \sigma \sinh \left( \| f \|_{L^\infty} \right) \leq 1 \, ,
\end{equation}
we have
\begin{equation}
\| \cF(u,v,w) \| \leq  R  \, ,
\end{equation}
that is, $\cF$ maps $\overline{B_R(0)}$ to itself.

Let us take $(u_1,v_1,w_1)$ and $(u_2,v_2,w_2)$ in $\overline{B_R(0)}$. We have
\begin{equation}
\| \bar{u}_1 - \bar{u}_2 \|_{L^\infty} \leq \Delta \tau \| v_1 - v_2 \|_{L^\infty} + \Delta \tau \| w_1 - w_2 \|_{L^\infty}
\end{equation}
and, for $\tau + \lambda \leq 1 + \tau_0$,
\begin{equation}
| \bar{v}_1 - \bar{v}_2 | \leq 2 \Delta \tau  \| K \|_{L^\infty} e^{2R} \| u_1 - u_2 \|_{L^\infty} \, .
\end{equation}
For $ \tau + \lambda\geq 1 + \tau_0$, on the other hand, we have
\begin{equation}
| \bar{v}_1 - \bar{v}_2 | \leq \| w_{\rig,1} - w_{\rig,2} \|_{L^\infty} + 2 \sigma \cosh R \| u_{\rig,1} - u_{\rig,2} \|_{L^\infty} + 2 \Delta \tau  \| K \|_{L^\infty} e^{2R} \| u_1 - u_2 \|_{L^\infty} \, .
\end{equation}
Since
\begin{equation}
\| w_{\rig,1} - w_{\rig,2} \|_{L^\infty} \leq 2 \Delta \tau  \| K \|_{L^\infty} e^{2R} \| u_1 - u_2 \|_{L^\infty} \, ,
\end{equation}
and
\begin{equation}
\| u_{\rig,1} - u_{\rig,2} \|_{L^\infty} \leq \Delta \tau \| v_1 - v_2 \|_{L^\infty} + \Delta \tau \| w_1 - w_2 \|_{L^\infty} \, ,
\end{equation}
we conclude that
\begin{equation}
\| \bar{v}_1 - \bar{v}_2 \|_{L^\infty} \leq 4 \Delta \tau  \| K \|_{L^\infty} e^{2R} \| u_1 - u_2 \|_{L^\infty} + 2\sigma \Delta \tau \cosh R \| v_1 - v_2 \|_{L^\infty} + 2\sigma \Delta \tau \cosh R \| w_1 - w_2 \|_{L^\infty}  \, ,
\end{equation}
and similarly
\begin{equation}
\| \bar{w}_1 - \bar{w}_2 \|_{L^\infty} \leq 4 \Delta \tau  \| K \|_{L^\infty} e^{2R} \| u_1 - u_2 \|_{L^\infty} + 2\sigma \Delta \tau \cosh R \| v_1 - v_2 \|_{L^\infty} + 2\sigma \Delta \tau \cosh R \| w_1 - w_2 \|_{L^\infty}  \, .
\end{equation}
In other words, we have
\begin{equation}
\| (\bar{u}_1,\bar{v}_1,\bar{w}_1) - (\bar{u}_2,\bar{v}_2,\bar{w}_2) \| \leq  M \Delta \tau \| (u_1, v_1, w_1) - (u_2, v_2, w_2) \|  \, ,
\end{equation}
where
\begin{equation} \label{sizeofM}
M = \max \left\{ 8 \| K \|_{L^\infty} e^{2R}, 1 + 4 \sigma \cosh R \right\}  \, .
\end{equation}
We conclude that $\cF$ is a contraction if $\Delta \tau$ is chosen small enough. Note from \eqref{sizeofR}, \eqref{sizeofdeltatau} and \eqref{sizeofM} that the choice of $\Delta \tau$ depends only on $\| f \|_{L^\infty}$, $\| g \|_{L^\infty}$, $\| f' \|_{L^\infty}$ and $\| K \|_{L^\infty}$.

The fixed point of $\cF$ yields a solution of the system of integral equations~\eqref{IBVP4} in the space of continuous functions. Since these equations imply
\begin{equation}
L^- v = L^+w  \, ,
\end{equation}
there exists a $C^1$ function $\tilde{u}$ such that\footnote{Recall that if $P(x,y)$ and $Q(x,y)$ are continuous and $\frac{\partial P}{\partial y} = \frac{\partial Q}{\partial x}$, where these partial derivatives exist and are continuous, then $\phi(x,y) = \int_0^x P(s,0) ds + \int_0^y Q(x,s) ds$ is a $C^1$ potential for $(P,Q)$.}
\begin{equation}
L^+\tilde{u} = v \, , \qquad L^-\tilde{u} = w \, ,
\end{equation}
and consequently
\begin{equation}
\dot{\tilde{u}} = L^+\tilde{u} + L^-\tilde{u} = v + w = \dot{u} \, ,
\end{equation}
that is,
\begin{equation} \label{uutilde}
\tilde{u}(\tau,\lambda) = u(\tau,\lambda) + h(\lambda)
\end{equation}
for some $C^1$ function $h$ (recall that $u(0,\lambda)=f(\lambda)$ is $C^2$). Since from~\eqref{IBVP4} we have
\begin{equation}
L^+u(\tau_0,\lambda) = f'(\lambda) + \frac12 v(\tau_0,\lambda) + \frac12 w(\tau_0,\lambda) = f'(\lambda) + g(\lambda) = v(\tau_0,\lambda) = L^+ \tilde{u}(\tau_0,\lambda) \, ,
\end{equation}
we obtain from \eqref{uutilde} that $h'(\lambda)=0$, that is, $u$ and $\tilde{u}$ differ by a constant, and so
\begin{equation}
L^+u = v \, , \qquad L^-u = w \, .
\end{equation}
Consequently, $u$ is in fact $C^1$, and so our hypotheses on the regularity of the initial data $f$ and $g$, together with \eqref{IBVP4}, show that $v$ and $w$ are also of class $C^1$ (note that this is not obvious along the curves $\tau + \lambda = 1 + \tau_0$ and $\tau - \lambda = \tau_0$, where the compatibility condition~\eqref{compatibility2} must be used). From this it is easily shown that $u$ is a $C^2$ solution of the IBVP~\eqref{IBVP3}.

Going back to the original IBVP~\eqref{IBVP2}, we see that its initial data satisfies the compatibility conditions~\eqref{compatibility1} and \eqref{compatibility2}. Consequently, it has a $C^2$ solution for $\tau \in [0, \Delta \tau]$. If we subsequently set $\tau_0 = \Delta \tau$, $f(\lambda)=u(\tau_0,\lambda)$ and $g(\lambda)=\dot{u}(\tau_0,\lambda)$, then $f$ and $g$ satisfy the compatibility conditions~\eqref{compatibility1} and \eqref{compatibility2}, and so the IBVP~\eqref{IBVP3} has a $C^2$ solution for $\tau \in [\tau_0, \tau_0 + \Delta \tau']$, meaning that the original IBVP~\eqref{IBVP2} has a $C^2$ solution for $\tau \in [0, \Delta \tau + \Delta \tau']$. Iterating this procedure, we can keep extending the solution to larger and larger time intervals. Since the time interval by which we can extend depends only on $\| f \|_{L^\infty}$, $\| g \|_{L^\infty}$, $\| f' \|_{L^\infty}$ and $\| K \|_{L^\infty}$, and the latter is smooth for $\tau \in [0, +\infty)$, we see that if $\| u(\tau, \cdot) \|_{L^\infty}$, $\| \dot{u}(\tau, \cdot) \|_{L^\infty}$ and $\| u'(\tau, \cdot) \|_{L^\infty}$ remain bounded then the solution is global in time. On the other hand, it is clear from the system of integral equations~\eqref{IBVP4} that $\| \dot{u}(\tau, \cdot) \|_{L^\infty}$ and $\| u'(\tau, \cdot) \|_{L^\infty}$ canot blow up if $\| u(\tau, \cdot) \|_{L^\infty}$ remains bounded. We conclude that either the solution is global or $\| u(\tau, \cdot) \|_{L^\infty}$ blows up.

Finally, we show that the solution is unique: given two $C^2$ solutions $u_1$ and $u_2$ defined on some compact time interval $[0, \tau_1]$, we take $R$ large enough such that both $(u_1,L^+{u}_1,L^-u_1)$ and $(u_2,L^+u_2,L^-u_2)$ are in the closed ball $\overline{B_R(0)}$ of the space $\left(C^0\left([0, \tau_1] \times [0,1]\right)\right)^3$, and choose $\Delta \tau = \tau_1 / N$ (for some $N \in \bbN$) small enough such that $\cF$ is a contraction on the closed ball $\overline{B_R(0)}$ of the space $\left(C^0\left([0,\Delta \tau] \times [0,1]\right)\right)^3$. Since both $(u_1,L^+{u}_1,L^-u_1)$ and $(u_2,L^+u_2,L^-u_2)$ are fixed points of $\cF$, they must coincide for $\tau \in [0,\Delta \tau]$. Iterating this procedure for time intervals of the form $[ k \Delta \tau, (k + 1) \Delta\tau]$, with $k = 1, \ldots, N-1$, we conclude that $u_1$ and $u_2$ coincide on $[0, \tau_1]$. 
\end{proof}

\end{document}
